\begin{afterword}
\summaryhead

Asked ``What is red?'', one can answer with more words (``a color,'' ``a property of a thing'') and soon be lost, like someone learning Chinese from a Chinese/Chinese dictionary. Or one can point to red things. The author calls the first an intensional definition and the second an extensional one, and recalls learning the distinction, perhaps from S. I. Hayakawa. Peirce's recipe for finding lithium is offered as a third kind: an intensional procedure that leads to an extensional example, a ``treasure map.''

The post then separates definitions from concepts. The actual intension of the author's ``tiger'' concept is the neural pattern that recognizes tigers; its actual extension is everything the author would call a tiger. Neither words nor pointing can convey the whole concept. So one cannot define a word any way one likes: concepts cannot be programmed into other people, and one cannot freely choose both a definition and what it picks out, as the example of Mars shows. Because this matching happens below awareness, arguments ``by definition'' can do work in debate that they could not do in reality.

\responsehead

The post teaches a real distinction, and its sources check. Harnad's paper does compare the grounding problem to learning Chinese from a Chinese/Chinese dictionary. The Peirce passage is quoted exactly, and Peirce's next sentence reads it as the post does, as ``a precept that is more serviceable than a definition.'' The hedged memory of Hayakawa is correct: \textsc{Language in Thought and Action} has the same regress about ``red'' and the same answer, the traffic light.

The terms, however, are Hayakawa's more than logic's, and a reader should know it. Hayakawa calls meaning given in more words ``intensional'' and meaning given by pointing ``extensional.'' In logic, an extensional definition lists every member of a class, and pointing to a few examples is an ostensive definition. The post's later exception, a concept whose entire extension it can show, uses the logicians' sense. Someone who takes the post's definitions into a logic textbook will find them used differently.

The step from definitions to concepts is the post's own, and it is a stipulation. ``The actual intension of my `tiger' concept'' is a neural pattern in one person's cortex, where logic's intension is a property that a word connotes. The post marks the shift, and its sense is close to the ``categorical representations'' of the Harnad paper it cites. Read in that sense, the post's claims about concepts are claims of psychology, and should not be taken as claims about what logicians mean by the word.

The two arguments that follow are sound. That pointing cannot fix which feature is meant is Wittgenstein's point about ostensive definition (1953). That a definition and its object must be matched separately is shown with Mars, whose figures are right. The last paragraph explains why arguments ``by definition'' are popular, but gives no example of one; later posts in the sequence do.

In short: a correctly sourced distinction, named in Hayakawa's sense rather than logic's, extended by a psychological sense of ``intension'' that the post stipulates, and ending on an explanation it does not illustrate.
\end{afterword}
